\documentclass[10pt]{article}
\usepackage{tmlr}
\usepackage{hyperref}
\usepackage{natbib}
\usepackage{booktabs}
\usepackage{graphicx}
\usepackage{subcaption}
\usepackage{amsmath, amssymb, amsthm}

\input{math_commands.tex}

\newtheorem{theorem}{Theorem}[section]
\newtheorem{proposition}[theorem]{Proposition}

% tighter spacing
\setlength{\floatsep}{3pt plus 1pt minus 1pt}
\setlength{\textfloatsep}{3pt plus 1pt minus 1pt}
\setlength{\intextsep}{4pt plus 1pt minus 1pt}
\setlength{\abovecaptionskip}{2pt}
\setlength{\belowcaptionskip}{2pt}
\setlength{\aboverulesep}{0pt}
\setlength{\belowrulesep}{0pt}
\addtolength{\belowcaptionskip}{-2pt}

% --- commands ---
\newcommand{\Kzero}{K_{0}}
\newcommand{\KT}{K_{T}}
\newcommand{\DK}{{\Vert K_{T}-K_{0}\Vert_{F}/\Vert K_{0}\Vert_{F}}}
\newcommand{\cca}{{\operatorname{CKA}}}
\newcommand{\frob}{{\operatorname{Frob}}}
\newcommand{\spearman}{{\rho_{s}}}

\title{When to Trust Kernel Change?\\A Predictive Validation Framework for Representation Diagnostics}

\author{
  \name Yaoping Wang \\
  \email yaoping.wang@example.com \\
  \addr Department of Statistics \\
  \addr University Name
}

\raggedbottom
\begin{document}

\maketitle

% ============================================================
\begin{abstract}
We establish finite-sample bounds (Proposition~\ref{prop:gap}) on the predictive validity of any label-invariant representation diagnostic, and propose RDEP, a principled validation framework that separates reliable diagnostics from misleading ones. Applying RDEP to feature-kernel evolution reveals two complementary findings. First, diagnostic reliability can be partially predicted from pre-training quantities alone: a regression across Hermite-family targets achieves $R^2=0.847$ (in-sample; leave-one-out CV $R^2=0.710$) using initial feature alignment and headroom, with headroom as the dominant univariate predictor (Spearman $\rho=0.81$, $p=0.0002$). Second, a task-aligned spectral correction that uses the same information \emph{within} a diagnostic fails under cross-validation ($\rho=-0.45$), exposing a deeper Diagnostic Validation Gap between in-sample and out-of-sample performance. Controlled experiments across multiple architectures (MLPs, CNNs, and Vision Transformers) and real datasets (FashionMNIST $\rho=-0.43$, CNN CIFAR-10) produce a reliability map, an adaptive diagnostic selection rule based on effective rank, and three validation principles applicable to any representation similarity metric.
\end{abstract}

% ============================================================
\section{Introduction}
\label{sec:intro}

\subsection{Background}

Representation similarity metrics (CKA, SVCCA, PWCCA, RSA, Frobenius distance) are routinely used to diagnose feature learning in neural networks. A high similarity between initial and trained representations is interpreted as evidence of limited learning; a large change is read as useful adaptation. These inferences influence how the field compares architectures, evaluates training dynamics, and judges representation quality. But they rest on an untested premise: that the magnitude of observable representation change is a sufficient statistic for generalization improvement.

The central claim of this paper is that \textbf{representation diagnostics themselves require predictive validation}. Just as we would not trust a model evaluated only on its training set, we should not trust a representation metric evaluated only on the same data where it was computed. We show that failing to validate leads to systematically misleading conclusions, and we describe a protocol for doing so.

This question connects directly to NTK theory~\citep{jacot2018neural}. The NTK predicts that infinitely wide networks converge to a fixed kernel, so kernel change should be near-zero at large widths. Our experiments confirm this prediction (CKA $>0.999$ at width~1024). The practical question for finite-width networks is not whether the kernel changes (NTK predicts it does not at infinite width) but \emph{when finite-width kernel change carries meaningful information about generalization}. We quantify the conditions under which kernel change is informative in finite-width regimes.

We introduce the \textbf{Representation Diagnostic Evaluation Methodology (RDEP)}. RDEP is a general framework for empirically validating whether a representation metric is informative about generalization and identifying when it fails (Figure~\ref{fig:reliability_map}). To our knowledge, it is the first application of predictive validation, a standard practice in supervised learning, to the evaluation of representation diagnostics. It bridges a gap between how we train models and how we assess their learned representations. This is not a paper about kernels. It is a paper about how the field evaluates its own diagnostic tools, and why that evaluation must be held to the same standard as the models being diagnosed. RDEP consists of four components: (i) a calibration table mapping task scale to metric--outcome correlation, (ii) width-controlled partial correlation analysis to remove spurious confounds, (iii) real-data validation across architectures, and (iv) held-out evaluation of task-aware diagnostics. We demonstrate RDEP on feature-kernel evolution (Frobenius, CKA, SVCCA, PWCCA) as a case study. We isolate the empirical feature kernel rather than the full NTK~\citep{jacot2018neural}, and produce a reliability map and validation principles that generalize across metric choices.

Across controlled experiments with three target functions, six network widths (32--1024), and per-epoch spectral tracking, we arrive at three empirical calibration points:

\begin{enumerate}
    \item \textbf{Kernel stability does not imply the absence of useful representation change.}  At large widths on the polynomial and high-frequency tasks, the kernel is virtually static (CKA $>0.999$, Frobenius change $\sim$1\%). Yet the network operates near its final representation with minimal room for improvement. Stability here reflects proximity to a fixed point, not failure to learn.  (GMM retains $\sim$12\% reorganization at width~1024, but still shows minimal improvement.)
    \item \textbf{Kernel evolution does not imply useful adaptation.}  At small widths, the kernel can reorganize substantially ($>100\%$ Frobenius change), but whether this translates into generalization improvement depends on the target: the per-dataset Spearman correlation ranges from $\rho=0.14$ (poly) to $\rho=0.94$ (highfreq) to $\rho=1.0$ (GMM).
    \item \textbf{CKA does not characterize spectral stability.}  Global similarity (CKA $>0.99$) can coexist with substantial spectral reorganization (Frobenius change $>40\%$), revealing that standard similarity metrics and spectral quantities carry distinct information.
\end{enumerate}

These findings form an empirical boundary: feature-kernel evolution magnitude alone is neither necessary nor sufficient as a diagnostic of finite-width generalization improvement.  The task relevance of specific representational changes, rather than their global magnitude, appears to be the essential quantity.  We advocate evaluating representation diagnostics by their ability to predict downstream learning behavior rather than solely by the magnitude of representation change they detect.

\paragraph{Contributions.}
We demonstrate RDEP on feature-kernel evolution (Frobenius, CKA, SVCCA, PWCCA) across synthetic and real data, producing (i) a calibration table mapping task structure and model scale to metric--outcome correlation, (ii) width-controlled partial correlation analysis unmasking spurious width-driven associations, (iii) validation principles for preventing data-reuse overfitting in task-aware diagnostics, and (iv) a practical four-step workflow for practitioners (detailed in Section~\ref{sec:principles}).

\paragraph{Three Core Principles.}
Our findings throughout the paper distill into three principles for evaluating representation diagnostics (elaborated in Section~5.3): (i) global similarity does not imply representation sufficiency, (ii) kernel change magnitude does not imply task relevance, and (iii) task alignment must be validated on held-out data.

\paragraph{Roadmap.}
Section~2 situates our work relative to NTK theory, spectral generalization bounds, and existing CKA critiques. Section~3 describes the experimental setup (models, data, metrics). Section~4 documents three reliability regimes where common diagnostic intuitions fail (Section~4.1--4.3) and validates the dissociation on real architectures (Section~4.4). Section~5 first asks whether diagnostic reliability can be predicted from pre-training quantities (Section~5.1.1 finds that it can; in-sample $R^2=0.864$, LOOCV $R^2=0.670$), then shows that the same spectral information cannot be naively turned into a cross-validated instance-level diagnostic (Section~5.2 reveals a negative correlation $\rho=-0.45$), and finally distills the findings into the three core principles above and a practical workflow (Section~5.3). Section~5.4 discusses limitations and Section~5.5 outlines future directions.

% ============================================================
\section{Related Work}
\label{sec:related}

\subsection{Neural Tangent Kernel and Infinite-Width Limits}
The NTK~\citep{jacot2018neural} connects infinitely wide neural networks to kernel methods. It provides tools for analyzing optimization~\citep{du2019neural,allen2019learning} and generalization~\citep{arora2019exact}.  Arora et al.~\cite{arora2019exact} characterized the convergence of empirical NTKs to their infinite-width limits.  The fixed-NTK description corresponds to the ``lazy training'' regime~\citep{chizat2020lazy}.  Our large-width experiments operate near this regime.  Our work does not challenge NTK theory within its intended regime; instead, we investigate whether observable feature-kernel evolution provides a sufficient explanation of finite-width feature learning across regimes.

\subsection{Feature Learning Beyond the NTK Regime}
Yang \& Hu~\cite{yang2021tensor} introduced maximal update parameterization ($\mu$P) to preserve feature learning in infinite-width limits.  Bordelon \& Pehlevan~\cite{bordelon2022spectral} and Lauditi et al.~\cite{lauditi2025adaptive} described feature-learning networks as kernel machines with data-dependent kernels, a framework in which kernel evolution is \emph{constitutive} of feature learning by definition.  Geiger et al.~\cite{geiger2020scaling} characterized the lazy-to-rich transition, and Kumar et al.~\cite{kumar2024grokking} used CKA to track this transition in grokking, treating kernel change as a diagnostic of feature learning rather than evaluating the diagnostic itself.  Bordelon \& Pehlevan~\cite{bordelon2023dynamics} analyzed finite-width kernel fluctuations, showing that feature evolution increases with deviation from the infinite-width limit.  Chou et al.~\cite{chou2025feature} further examined representational geometry across the lazy-to-rich spectrum.  These works primarily address the \emph{existence} and \emph{character} of kernel evolution during training, often using CKA or similar metrics as measurement tools.  If kernel evolution is constitutive of feature learning, as the adaptive kernel framework suggests, then the relevant question is not \emph{whether} the kernel changes but \emph{when} that change is informative about generalization.  Our work examines this question empirically: we measure how often and under what conditions observable kernel evolution provides a sufficient diagnostic for generalization improvement.

\subsection{Spectral Analysis of Kernel Generalization}
Canatar, Bordelon, and Pehlevan~\cite{canatar2021spectral} derived the generalization error of kernel regression in terms of eigenvalues $\lambda_{i}$ and target projections $v_{i}=\langle f^{*},\phi_{i}\rangle^{2}$.  This motivates our methodology: rather than global similarity metrics, we examine spectral quantities that directly enter kernel risk estimates.

\subsection{Representation Similarity Metrics}
CKA~\citep{kornblith2019similarity} is widely used for representation comparison~\citep{raghu2017svcca,morcos2018insights,klabunde2023similarity}.  CKA measures global alignment of kernel matrices.  Our results show that global alignment and spectral stability are distinct quantities.  High CKA does not imply spectral stability, a distinction that matters when kernel analyses depend on spectral quantities.

Existing critiques of CKA~\citep{davari2022reliability,ding2021grounding} primarily focus on statistical reliability, asking whether CKA values are stable across samples, sensitive to dimensionality, or biased by network width.  Kang et al.~\cite{kang2025spectral} further examined how limited neuron sampling biases CKA, proposing a spectral denoising correction.  Dujmović et al.~\cite{dujmovic2025pitfalls} documented analogous pitfalls for RSA.  Murphy et al.~\cite{murphy2024cka} and Chun et al.~\cite{chun2025debias} developed finite-sample debiased CKA estimators, directly relevant to our discussion of the effective rank bound $p_{\text{eff}} / n$ in Proposition~\ref{prop:gap}.  These are important concerns, but they address a different question from ours: they ask whether CKA accurately measures what it claims to measure, while we ask whether that measured quantity is sufficient for predicting generalization.  Even if CKA were perfectly stable and unbiased, the magnitude of kernel change it measures would still be insufficient for diagnosing generalization, because the generalization error depends on the spectral direction of change, not its global magnitude.  Our contribution is to isolate this diagnostic insufficiency empirically, across a calibrated range of tasks and widths, rather than to critique the metric's statistical or numerical properties.

\begin{table}[h]
\centering\small
\caption{How this work differs from prior CKA critiques.}
\label{tab:related_diff}
\begin{tabular}{p{3cm}p{3.2cm}p{3.2cm}p{3.2cm}}
\toprule
Dimension & Davari et al.~\cite{davari2022reliability} & Ding et al.~\cite{ding2021grounding} & This work \\
\midrule
Core question & Is CKA statistically stable? & Is CKA biased by width? & Is kernel change sufficient for predicting generalization? \\
Validation required & Bootstrap resampling & Width normalization & Held-out predictive validation (RDEP) \\
Evaluation scope & Single metric (CKA) & Single metric (CKA) & Five metrics (CKA, SVCCA, PWCCA, RSA, Frobenius) \\
Task structure & Not varied & Not varied & Varied (3 targets, 6 widths) \\
\bottomrule
\end{tabular}
\end{table}

In short, prior work asks ``is the metric stable?'' while we ask ``is the metric sufficient?'' These are complementary but distinct questions, and answering the first does not answer the second.  Concurrent work by~\citet{bridging2026usable} argues that representational similarity is sufficient but not necessary for functional similarity, reinforcing our findings from a complementary perspective.

\subsection{Theoretical Motivation: The Diagnostic Validation Gap}
\label{sec:proposition}

Any label-invariant global similarity diagnostic $\mathcal{D}(K_0, K_T)$ (CKA, SVCCA, PWCCA, RSA, Frobenius distance) depends only on the kernel matrices and not on task labels $y$. The generalization error $\mathcal{E}$ depends on the interaction between spectral quantities and the target through the projections $v_i$~\citep{canatar2021spectral}.  Consequently, $\mathcal{D}$ and $\Delta\mathcal{E}$ can in principle dissociate.  While the possibility of dissociation follows from basic spectral considerations, two questions remain open: first, whether a diagnostic that \emph{incorporates} target information via estimated projections $\hat{v}_i$ suffers from a systematic validation gap when those projections are estimated from finite data; and second, how prevalent the dissociation is across representative learning regimes.  The first question is a novel statistical phenomenon---the \textbf{Diagnostic Validation Gap}---which we formalize below.  The second we address empirically through the RDEP framework.

We formalize the dissociation as a finite-sample phenomenon using the spectral risk framework of~\citet{canatar2021spectral}.  Let $K_0 = \sum_{i=1}^n \lambda_i \phi_i \phi_i^\top$ be the eigendecomposition of the initial kernel, and let $v_i = \langle \phi_i, y \rangle^2 / n$ be the target projections.  Consider a task-aligned diagnostic
\[
\mathcal{D}_{\text{align}} = \sum_{i=1}^n \hat{v}_i \cdot \Delta\lambda_i,
\]
where $\hat{v}_i$ are estimated from a training sample and $\Delta\lambda_i = \lambda_i^{(T)} - \lambda_i^{(0)}$ are the eigenvalue shifts after training.  The generalization improvement $\Delta\mathcal{E}$ also depends on $\{v_i\}$ through the spectral risk estimate, creating a statistical coupling that underlies the validation gap.

\begin{proposition}[Diagnostic Validation Gap]
\label{prop:gap}
Let $\hat{v}_i$ be estimated from $n$ samples and let $v_i^*$ be the population target projections.  Let $\rho_{\text{in}}$ be the Spearman correlation between $\mathcal{D}_{\text{align}}$ and $\Delta\mathcal{E}$ when both are computed on the same $n$ samples, and $\rho_{\text{cv}}$ the correlation under $K$-fold cross-validation where $\hat{v}_i$ and $\Delta\mathcal{E}$ are measured on independent splits.  Then
\[
\mathbb{E}[\rho_{\text{in}}] = \rho^* + \mathcal{O}\bigl(\|\hat{v} - v^*\|_1\bigr),\qquad
\operatorname{Var}(\rho_{\text{cv}}) = \Omega\!\left(\frac{p_{\text{eff}}}{n}\right),
\]
where $\rho^*$ is the oracle correlation that would be obtained with known $v_i^*$, $p_{\text{eff}} = (\sum \lambda_i)^2 / \sum \lambda_i^2$ is the effective rank of $K_0$, and $\Omega$ denotes asymptotic lower bound.  The in-sample estimate carries a systematic optimism bias $\mathcal{O}(\|\hat{v} - v^*\|_1)$ due to double-use of the training labels.  Cross-validation removes this bias but introduces variance that scales with $p_{\text{eff}} / n$, which can dominate when the effective rank is comparable to the sample size.
\end{proposition}

The proof follows from two standard results. The in-sample bias $\mathcal{O}(\|\hat{v} - v^*\|_1)$ arises because $\hat{v}_i$ and the residual $y - \hat{y}$ are correlated through the shared training labels, a finite-sample analogue of the optimism of the training error~\citep{efron1983estimating}. The cross-validation variance $\Omega(p_{\text{eff}} / n)$ follows from the observation that $\mathcal{D}_{\text{align}} = \sum_i \hat{v}_i \Delta\lambda_i$ is a weighted sum of $p_{\text{eff}}$ effectively independent terms when the kernel spectrum is flat; under $K$-fold cross-validation, the estimates $\hat{v}_i$ from different folds are uncorrelated, leading to a variance lower bound that scales as $p_{\text{eff}} / n$. A detailed derivation is provided in Appendix~A.

Proposition~\ref{prop:gap} explains our central negative result ($\rho = -0.45$ under cross-validation in Section~5.2): with $n=240$ training samples and a feature kernel whose effective rank at large widths exceeds $n$, the cross-validated variance term $\Omega(p_{\text{eff}} / n)$ is large enough to render the estimate unreliable---and in our experiment, to flip its sign.  This is not a failure of the task-aligned idea but a fundamental statistical limitation: the data used to estimate $v_i$ cannot simultaneously serve to validate the diagnostic that uses them.

Proposition~\ref{prop:gap} also suggests a potential remedy: the variance of $\rho_{\text{cv}}$ scales with $p_{\text{eff}} / n$, implying that regularization targeted at reducing the effective rank could stabilize the estimate.  A ridge-style shrinkage estimator $\tilde{v}_i = \lambda_i / (\lambda_i + \gamma) \cdot \hat{v}_i$ would in principle reduce variance at the cost of bias, with the optimal trade-off governed by $p_{\text{eff}} / n$.  However, in our experimental setting the initial feature kernel has a near-uniform eigenvalue spectrum (all $\lambda_i$ are of similar magnitude), so the shrinkage factor is approximately constant across all directions and provides no discriminatory benefit.  We verified this empirically: applying shrinkage with a range of $\alpha$ values left the cross-validated correlation essentially unchanged ($\rho \approx -0.66$ for all $\alpha$ at width~1024).  Shrinkage-based corrections may be more effective in settings with clearer spectral decay, such as the NTK regime or large-scale pre-trained models, which we leave to future work.

% ============================================================
\section{Experimental Setup}
\label{sec:setup}

\subsection{Models}
Two-layer ReLU networks (widths 32--1024, no bias) with He initialization:
\[
f(x)=W_{2}\cdot\operatorname{ReLU}(W_{1}x),\quad
W_{1}\in\mathbb{R}^{m\times d},\; W_{2}\in\mathbb{R}^{1\times m},
\]
where $m$ is the width.  Standard parameterization (not $\mu$P/NTK parameterization).  SGD with momentum $0.9$, learning rate $0.05$, no weight decay, $500$ epochs, MSE loss, gradient clipping at $1.0$.

\subsection{Data}
$n=300$, $d=10$, 80/20 train/val split.  Three target functions:
\begin{description}
    \item[poly:] $f(x)=x_{1}^{2}+x_{2}$, inputs $\sim\operatorname{Uniform}[-1,1]^{d}$.
    \item[highfreq:] $f(x)=\sin(5x_{1})+\cos(7x_{2})$, inputs $\sim\operatorname{Uniform}[-1,1]^{d}$.
    \item[gmm:] $y\in\{0,1\}$ for two isotropic Gaussian clusters separated by $2$ units.
\end{description}
Targets standardized to zero mean and unit variance.  Label noise $\sigma=0.1$.

\subsection{Feature Kernel}
We define the \textbf{empirical feature kernel} (representation Gram matrix):
\[
K(x,y)=\frac{1}{m}\,h(x)^{\top}h(y),\quad h(x)=\operatorname{ReLU}(W_{1}x).
\]
Its eigendecomposition $K=U\Lambda U^{\top}$ gives eigenvalues $\Lambda=\{\lambda_{i}\}$ and eigenvectors $U=\{\phi_{i}\}$. These describe the sample-space geometry induced by the learned representations.  We focus on this quantity rather than the full NTK to isolate representation evolution from optimization dynamics.  All kernel comparisons are performed on the fixed evaluation set (all $n$ points), so changes reflect representation evolution rather than changes in sampling geometry.

\subsection{Metrics}
\begin{itemize}
    \item \textbf{CKA}~\citep{kornblith2019similarity}: centered kernel alignment between initial and trained kernels. Intuitively, CKA measures how similar two representation similarity matrices are as a whole (like comparing two images by overall layout). This makes it insensitive to changes in individual eigen-directions.
    \item \textbf{Frobenius change:} $\DK$. This measures the total per-entry magnitude of kernel change (like the sum of pixel-wise differences between two images). It captures changes across the full spectrum, including directions where the target has little projection.
    \item \textbf{Eigenvalue ratio:} $\lambda_{T}(i)/\lambda_{0}(i)$ per eigen-direction.
    \item \textbf{Target projection:} $v_{i}=\langle y,\phi_{i}\rangle^{2}/n$~\citep{canatar2021spectral}.
    \item \textbf{Generalization improvement:} $-\Delta\text{Error}=\text{Error}_{0}-\text{Error}_{T}$.
\end{itemize}

\subsection{Statistical Reporting}
Synthetic experiments use 10 random seeds per configuration (5 for real-data experiments).  Reported values are means $\pm 1$ standard deviation across seeds.  Spearman correlations are computed across configuration means (18 points: 3 datasets $\times$ 6 widths, or 6 points within each dataset) to avoid pseudo-replication from the strong width-driven trend within each condition.  Bootstrap 95\% confidence intervals are reported where noted.

\subsection{Hyperparameters}
All experiments use the following settings unless noted otherwise.  \textbf{Architecture:} two-layer ReLU MLP with hidden widths 32--1024, no bias, He-like initialization (weights $\sim \mathcal{N}(0, \sqrt{2/d_{\text{in}}})$).  \textbf{Optimization:} SGD with momentum 0.9, learning rate 0.05 (no schedule), no weight decay, 500 epochs, MSE loss, gradient clipping at 1.0.  \textbf{Data:} $n=300$ samples, $d=10$ input dimensions, 80/20 train/validation split, label noise $\sigma=0.1$.  \textbf{CNN (CIFAR-10):} two convolutional layers (filter size $5\times5$, channels $W$ and $2W$ for $W\in\{16,32,64,128,256\}$), each followed by ReLU and $2\times2$ max-pooling, then a fully-connected layer to 1 output.  SGD with momentum 0.9, learning rate 0.01, 100 epochs.  \textbf{MNIST/FashionMNIST:} same MLP architecture as synthetic experiments, with 784 input units, learning rate 0.05, 500 epochs.

% ============================================================
\section{Calibrating Diagnostic Reliability}
\label{sec:results}

\subsection{Reliability Regime I: Stability without Learning}
\label{sec:stability}

\begin{figure}[t]
    \centering
    \includegraphics[width=0.8\textwidth]{fig1.png}
    \caption{
        \textbf{Width controls feature-kernel stability.}
        (a) CKA$(K_{0},K_{T})$ as a function of width.
        At width 1024, CKA exceeds 0.999 across all tasks.
        At width 32, CKA drops to 0.85--0.99 depending on the target.
        (b) Normalized Frobenius change $\|K_{T}-K_{0}\|_{F}/\|K_{0}\|_{F}$.
        At width 1024, poly and highfreq drop to $\sim$1\%.
        At width 32, it reaches 25--103\% depending on the target.
        Error bars show $\pm1$ s.d. over 10 seeds.
    }
    \label{fig:stability}
\end{figure}

\noindent\textbf{Intuition says:} A stable kernel (CKA $>0.99$) means the network did not learn anything new.

\noindent\textbf{What we found:} At width~1024, CKA exceeds $0.999$, but the network was already near-optimal from initialization. Stability did not mean ``no learning''; it meant ``no room to learn.''

Width controls feature-kernel stability (Figure~\ref{fig:stability}).  For the polynomial and high-frequency targets, CKA rises above $0.999$ and Frobenius change drops to $\sim$1\% at width~1024.  At width~32, CKA drops to $0.85$--$0.92$ and Frobenius change reaches $25$--$103\%$ depending on the target.  The same stability trend is observed on MNIST (Table~\ref{tab:mnist}): CKA ranges from $0.77$ (width~32) to $0.96$ (width~256), though with a non-monotonic deviation at width~1024 (see Table~\ref{tab:mnist} caption).  The GMM task shows a distinct pattern: even at width~1024, Frobenius change remains $\sim$12\%. This suggests that kernel evolution on certain easy-separation tasks continues beyond the width-driven stabilization observed for other targets.

The relationship between stability and learning is not straightforward.  For the polynomial target, the average generalization improvement is near-zero across \emph{all} widths ($\Delta$Error $0.003$--$0.023$; Table~\ref{tab:results}). Yet the kernel undergoes substantially different reorganization, from $\sim$25\% at width~32 to $\sim$1\% at width~1024.  Despite similar outcomes, the degree of kernel evolution varies by an order of magnitude: kernel change magnitude alone is not a reliable indicator of whether learning has occurred.  Regimes of maximal kernel stability are regimes where training-driven improvement is minimal, but the converse (that large kernel change signals useful adaptation) is not generally true.

Training dynamics for poly width 32 (Supplementary Figure~S4) reinforce this conclusion. CKA stabilizes within 50 epochs while most of the 500-epoch improvement occurs afterward. However, this improvement is seed-dependent: across 10 seeds, the mean improvement at width~32 is $0.003\pm0.049$, so the trajectory shown illustrates the temporal decoupling mechanism rather than a representative effect size.

\begin{table}[t]
    \centering\footnotesize
    \setlength{\tabcolsep}{4pt}
    \caption{Calibration table for feature-kernel diagnostics (10 seeds per configuration).  Values are means $\pm$ s.d.  CKA, SVCCA, PWCCA, $\Delta K$, and $\Delta$Error are reported for each (dataset, width) pair.  All three similarity metrics show the same qualitative pattern: high values coexist with a wide range of $\Delta$Error.}
    \label{tab:results}
    \begin{tabular}{lcccccc}
    \toprule
    Dataset & Width & CKA & SVCCA & PWCCA & $\Delta K$ (\%) & $\Delta$Error \\
    \midrule
    poly & 32   & $0.916\pm0.021$ & $0.867\pm0.020$ & $0.911\pm0.011$ & $25.2\pm3.0$ & $+0.003\pm0.049$ \\
    poly & 64   & $0.977\pm0.004$ & $0.932\pm0.006$ & $0.958\pm0.006$ & $12.5\pm1.0$ & $+0.023\pm0.030$ \\
    poly & 128  & $0.996\pm0.001$ & $0.985\pm0.002$ & $0.986\pm0.004$ & $5.2\pm0.5$  & $+0.013\pm0.015$ \\
    poly & 256  & $0.999\pm0.000$ & $0.992\pm0.004$ & $0.997\pm0.001$ & $3.0\pm1.3$  & $+0.008\pm0.004$ \\
    poly & 512  & $1.000\pm0.000$ & $0.999\pm0.001$ & $1.000\pm0.000$ & $2.3\pm0.8$  & $+0.006\pm0.002$ \\
    poly & 1024 & $1.000\pm0.000$ & $1.000\pm0.000$ & $1.000\pm0.000$ & $1.1\pm0.4$  & $+0.003\pm0.001$ \\
    \midrule
    highfreq & 32  & $0.851\pm0.020$ & $0.777\pm0.014$ & $0.779\pm0.041$ & $102.9\pm14.2$ & $+0.190\pm0.130$ \\
    highfreq & 64  & $0.942\pm0.009$ & $0.866\pm0.006$ & $0.857\pm0.033$ & $50.9\pm5.6$   & $+0.039\pm0.036$ \\
    highfreq & 128 & $0.985\pm0.002$ & $0.953\pm0.005$ & $0.945\pm0.008$ & $20.4\pm2.0$   & $+0.008\pm0.014$ \\
    highfreq & 256 & $0.997\pm0.000$ & $0.969\pm0.004$ & $0.967\pm0.007$ & $6.5\pm1.3$    & $+0.003\pm0.003$ \\
    highfreq & 512 & $1.000\pm0.000$ & $0.980\pm0.005$ & $0.976\pm0.006$ & $2.4\pm1.5$    & $-0.000\pm0.002$ \\
    highfreq & 1024 & $1.000\pm0.000$ & $0.994\pm0.002$ & $0.989\pm0.003$ & $0.8\pm0.2$    & $+0.000\pm0.000$ \\
    \midrule
    gmm & 32  & $0.992\pm0.006$ & $0.943\pm0.035$ & $0.984\pm0.014$ & $35.7\pm15.7$ & $+0.158\pm0.104$ \\
    gmm & 64  & $0.993\pm0.005$ & $0.976\pm0.011$ & $0.995\pm0.003$ & $41.5\pm20.3$ & $+0.228\pm0.150$ \\
    gmm & 128 & $0.985\pm0.021$ & $0.977\pm0.022$ & $0.993\pm0.012$ & $53.3\pm10.4$ & $+0.320\pm0.131$ \\
    gmm & 256 & $0.997\pm0.003$ & $0.990\pm0.008$ & $0.998\pm0.002$ & $34.4\pm8.4$  & $+0.144\pm0.079$ \\
    gmm & 512 & $0.999\pm0.002$ & $0.994\pm0.007$ & $0.998\pm0.003$ & $22.2\pm8.6$  & $+0.069\pm0.051$ \\
    gmm & 1024 & $1.000\pm0.000$ & $1.000\pm0.000$ & $1.000\pm0.000$ & $11.7\pm4.8$  & $+0.023\pm0.020$ \\
    \bottomrule
    \end{tabular}
\end{table}

\subsection{Reliability Regime II: Change without Predictable Learning}
\label{sec:change}

\begin{figure}[t]
    \centering
    \includegraphics[width=\textwidth]{fig2.png}
    \caption{
        \textbf{Similarity metrics, generalization improvement, and the CKA blind spot.}
        (a--c) CKA, SVCCA, and PWCCA versus generalization improvement ($-\Delta$Error), colored by dataset.
        All three metrics show the same qualitative pattern: high metric values coexist with a wide range of generalization outcomes.
        (d) CKA versus Frobenius change.  The shaded region (CKA $>0.99$) spans Frobenius changes from $\sim1\%$ to $>40\%$, confirming the CKA blind spot.
        Per-dataset Spearman correlations are reported in each panel.
    }
    \label{fig:dissociation}
\end{figure}

\noindent\textbf{Intuition says:} A large kernel change ($>100\%$ Frobenius) means the network learned something useful.

\noindent\textbf{What we found:} The correlation between kernel change and generalization improvement ranges from $\rho=0.14$ (poly) to $\rho=1.0$ (GMM), but the $\rho=0.94$ correlation for highfreq collapses to $\rho=0.10$ after controlling for width. Large change is often just noise.

If kernel evolution were a reliable indicator of learning progress, its magnitude should correlate with generalization improvement.  Pooled across all 18 configuration means, the Spearman correlation between kernel change (Frobenius \%) and generalization improvement is $\rho=0.83$ (95\% CI $[0.53, 0.96]$, $p=2\times10^{-5}$).  However, this aggregate association is almost entirely driven by between-task differences in average kernel change and improvement. Within each dataset, the relationship varies dramatically: $\rho=0.14$ (polynomial) to $\rho=0.94$ (high-frequency) and $\rho=1.0$ (GMM).  We caution that the poly estimate is based on only 6 configuration means (95\% CI spans $[-0.80, 1.00]$). To quantify how much of this uncertainty reflects limited width resolution, we ran an extended sweep with 11 widths (adding 5 intermediate values). The poly estimate becomes $\rho=0.36$ (95\% CI $[-0.43, 0.97]$, $p=0.27$), confirming that the null result is not an artifact of low resolution but a genuine statistical limitation. The wider grid does not change the qualitative conclusion: kernel change magnitude provides no statistically detectable signal for this task. With this caveat, we examine each task separately.

Within-dataset patterns reveal the dissociation more clearly:
\begin{description}
        \item[poly ($\rho=0.14$, 95\% CI ${[-0.80, 1.00]}$, $p=0.79$; extended sweep $\rho=0.36$, 95\% CI ${[-0.43, 0.97]}$, $p=0.27$):] No statistically detectable relationship (CI spans zero).  Despite $\sim25\%$ kernel reorganization at width~32, the average improvement is $0.003\pm0.049$, indistinguishable from zero.  Kernel change magnitude alone carries no predictive signal for this task. The wide confidence interval reflects the limited statistical power of 6 configuration means, but expanding to 11 widths narrows the CI without changing the qualitative conclusion.
    \item[highfreq ($\rho=0.94$, 95\% CI ${[0.52, 1.00]}$, $p=0.005$):] Strong positive correlation driven primarily by width variation: kernel change and generalization improvement both decline monotonically with increasing width.  At width~32, mean kernel change reaches $102.9\%$ with gain $0.190\pm0.130$ (high seed-to-seed variability).  While this width-driven relationship is clear, the correlation across widths does not translate into reliable prediction for any fixed width. Within width~32, individual seeds span a wide range of outcomes despite equivalent kernel change.
    \item[gmm ($\rho=1.0$, 95\% CI ${[1.00, 1.00]}$, $p<0.001$):] Near-perfect positive correlation, contrasting sharply with the poly and highfreq results.  GMM is an easy-separation task where small-width networks improve substantially, and kernel change tracks this improvement nearly monotonically.  The contrast among the three datasets demonstrates that kernel change magnitude alone is not a task-invariant diagnostic: its informativeness is itself task-dependent.

\end{description}

The same task-dependent pattern holds across all five metrics we examined, including RSA (Representational Similarity Analysis), which measures Spearman correlation between Euclidean distance matrices and is mathematically distinct from kernel alignment. For the polynomial task, all similarity metrics show no detectable correlation with $\Delta$Error (RSA $\rho=-0.43$, SVCCA $\rho=-0.43$, PWCCA $\rho=-0.43$, CKA $\rho=-0.43$, all $p>0.3$), consistent with the Frobenius result showing no detectable correlation. For highfreq, all similarity metrics show perfect negative correlation (RSA $\rho=-1.00$, $p<0.001$), simply reflecting that higher similarity means less reorganization rather than more learning. For GMM, the metrics show moderate negative correlation (RSA $\rho=-0.60$, SVCCA $\rho=-0.49$, PWCCA $\rho=-0.60$) while Frobenius remains near-perfect ($\rho=1.00$). This consistency across metrics, including one based on an entirely different mathematical framework (distance correlation rather than kernel alignment), confirms that the dissociation is a general property of label-invariant similarity diagnostics when evaluated without task-aware validation.

\subsubsection{Width-Controlled Analysis}

The per-dataset correlations above pool across widths. This raises a question: are they driven by genuine task structure or simply by the width-driven trend (wider networks change less and improve less)? To disentangle these, we compute partial Spearman correlations controlling for width and examine fixed-width seed-level variation.

For highfreq, the raw correlation $\rho=0.94$ collapses to $\rho=0.10$ after controlling for width, and the mean within-width correlation across seeds is $\rho=0.03$. This is consistent with the association being an artifact of the width-driven trend rather than a genuine relationship between kernel reorganization and generalization. For poly, controlling for width leaves the correlation near-zero ($\rho=0.08$). For GMM, the correlation persists at $\rho=0.91$ even within fixed widths, indicating a genuine task-level relationship between seed-level kernel variation and improvement. These results show that the relationship between kernel change and generalization depends on the task, and the width confound alone cannot explain the full pattern. We caution that these partial correlations are computed from 6 configuration means and should be interpreted as suggestive rather than definitive; bootstrap confidence intervals span the full range of possible values due to the small sample.

On MNIST binary classification, the Spearman correlation between kernel change and generalization improvement is $\rho=0.66$ (95\% CI $[0.36, 0.85]$, $p<0.001$; Supplementary Figure~S2), higher than on synthetic data but still far from a reliable sufficient diagnostic.  The stronger association likely reflects the simpler structure of the MNIST task, where initial representations are further from optimal and leave more room for task-aligned kernel adaptation.  Even here, kernel change magnitude alone does not uniquely determine generalization gains.

\textbf{Random-label baseline: kernel change under the null.}  To establish how much kernel change occurs in the absence of any signal, we trained the same architecture on randomly shuffled labels for all three datasets.  Under the null, kernel change remains substantial at small widths across all tasks (e.g., highfreq: $\sim\!108\%$ Frobenius at width~32; GMM: $\sim\!47\%$) and decays with width.  These magnitudes are comparable to (and at small widths exceed) those observed with structured targets. This places a strong lower bound on the claim that Frobenius change is task-relevant: the metric cannot even satisfy the \emph{necessary} condition of being larger in the presence of a signal than in its absence.  At width~32, Frobenius change under shuffled labels ($103.0\%$) exceeds that observed with the structured polynomial target ($25.2\%$), implying that a large fraction of the kernel reorganization observed in practice may reflect optimization dynamics unrelated to the task.

The null Spearman correlations between kernel change and $\Delta$Error are $\rho=0.53$ (poly, $p=0.002$), $\rho=0.66$ (highfreq, $p<0.001$), and $\rho=0.30$ (GMM, $p=0.11$).  The highfreq null correlation is larger than the real poly signal ($\rho=0.14$). This shows that the metric--outcome correlation statistic itself is noise-dominated and can appear informative even without a true task.  These null correlations are attributable to the strong width-dependence of both quantities rather than any meaningful association.

\subsection{Reliability Regime III: The CKA Blind Spot}
\label{sec:cka}

\noindent\textbf{Intuition says:} CKA $>0.99$ means the kernel is spectrally stable.

\noindent\textbf{What we found:} CKA $>0.99$ coexists with Frobenius changes $>40\%$ and substantial eigenvalue reorganization. Global similarity and spectral stability are distinct quantities.

CKA between initial and final kernels is high across nearly all configurations (Figure~\ref{fig:dissociation}d).  However, CKA and Frobenius change are not redundant.  Configurations with $\cca>0.99$ span a Frobenius range from $\sim1\%$ to $>40\%$.  The GMM case is the most striking: $\cca$ remains above $0.985$ across all widths, yet Frobenius change ranges widely (from $\sim$10\% to $\sim$50\%).

This distinction matters because the quantities that enter kernel risk estimates~\citep{canatar2021spectral} are spectral, namely eigenvalues $\lambda_{i}$ and target projections $v_{i}=\langle y,\phi_{i}\rangle^{2}/n$, not global matrix similarities.  These can undergo substantial reorganization that CKA averages out.  From a random matrix perspective, this dissociation is natural. CKA is dominated by alignment of the leading eigenspaces. Frobenius change captures the entire spectrum including the bulk, which undergoes fluctuation that a rank-sensitive metric like CKA does not see~\citep{el1998concentration,vyas2023statistical}.

In GMM at width 128, for instance, the second eigenvalue drops to $0.62\times$ its initial value while $v_{i}$ shifts by up to $4.4\times$ across eigen-directions, both invisible to CKA.  Supplementary Figure~S1 shows the full spectral energy redistribution across datasets.

The GMM case is particularly instructive because it exhibits the strongest association between kernel change and generalization improvement ($\rho=1.0$).  Unlike the polynomial task, where kernel reorganization at small widths yields near-zero average improvement, GMM's kernel changes at small widths are accompanied by large and consistent generalization gains (e.g., $\Delta$Error $=0.32\pm0.13$ at width~128).  The near-perfect correlation suggests that when the target function is well-aligned with the network's inductive bias (a binary classification task in two well-separated clusters), kernel magnitude evolution can closely track learning progress.  The contrast among the three tasks shows that even this strong association is task-dependent. The same diagnostic (kernel change magnitude) is highly informative for GMM, uninformative for poly, and noise-dominated for highfreq. All under identical experimental conditions.

\subsection{Reliability on Real Data}
\label{sec:real}

To test whether the dissociation extends beyond synthetic 10D tasks, we evaluate the feature-kernel diagnostic on two real-image benchmarks: CNN on CIFAR-10 and FashionMNIST.

\begin{figure}[t]
    \centering
    \includegraphics[width=\textwidth]{fig4.png}
    \caption{\textbf{Real-data validation.}  (a) CNN on CIFAR-10: Frobenius change versus linear probe $\Delta R^2$ (5 widths, 5 seeds).  The CKA blind spot ($0.991$ CKA with $43\%$ Frobenius) persists for deeper architectures on real images.  (b) FashionMNIST binary: kernel change versus $\Delta$Error.  Larger kernel changes associate with \emph{less} improvement ($\rho=-0.43$).}
    \label{fig:real_data}
\end{figure}

CNN on CIFAR-10 (2 conv layers + FC, channel widths $W=16,32,64,128,256$; Figure~\ref{fig:real_data}a) shows a strong positive Spearman correlation between Frobenius change and linear probe $\Delta R^2$ ($\rho=0.83$, $p<0.001$). On real images with deeper architectures, larger kernel change does track larger representation improvement. The dissociation is less extreme than on synthetic tasks. However, the relationship is non-monotonic: improvement peaks at $W=32$ ($+0.52$) and declines at both smaller and larger widths, while Frobenius change follows a different pattern. At $W=128$, CKA reaches $0.991$ with $43\%$ Frobenius reorganization. The CKA blind spot persists even though the Frobenius--$\Delta R^2$ relationship is generally positive. The dissociation is architecture- and data-dependent: real-world tasks can exhibit stronger association than synthetic benchmarks. The implication is clear: per-case calibration matters more than blanket claims.

FashionMNIST (Figure~\ref{fig:real_data}b) shows a negative Spearman correlation between kernel change and improvement ($\rho=-0.43$, $p=0.017$): larger kernel changes are associated with less improvement, and in some configurations performance degrades after training.  Kernel change remains high across all widths ($>67\%$ Frobenius).  Training dynamics at width~64 (5 seeds) reveals the mechanism: training MSE drops to near-zero ($0.0001\pm0.0002$) while test MSE remains substantial ($0.12\pm0.03$). The observed kernel reorganization reflects overfitting to nuisance variation rather than task-aligned adaptation.  Together, these results show that the diagnostic failures documented on synthetic tasks are not artifacts of two-layer networks or low-dimensional data. On FashionMNIST, the diagnostic is not merely uninformative but actively misleading.

To test whether the dissociation extends to attention-based architectures, we trained a small Vision Transformer (ViT-tiny: 4 layers, 4 heads, embedding dimensions 64--256) on binary MNIST (classes 0 vs.\ 1). We emphasize that this is a proof-of-concept on a small-scale attention architecture, not a claim about large-scale Transformers; the purpose is solely to test whether the dissociation is an artifact of the MLP parameterization. The pattern mirrors the synthetic experiments: Frobenius change correlates positively with $\Delta$Error across embedding dimensions ($\rho=0.90$, $p=0.04$), driven by the width-like trend where larger embeddings undergo less reorganization and also improve less. Meanwhile, CKA remains near-constant ($0.54$--$0.64$) across all dimensions, showing no detectable relationship with improvement ($\rho=0.10$, $p=0.87$). The CKA blind spot and the width-driven dissociation are not artifacts of the MLP parameterization---they appear in attention-based architectures as well.

To test the headroom framework across modalities (Section~\ref{sec:discussion}), we additionally trained a small Transformer (4 layers, 4 heads, embedding dimensions 64--256) from scratch on the SST-2 sentiment classification task (binary, 3000 training samples). This provides a cross-regime test: unlike the synthetic experiments where headroom varies systematically with task structure, SST-2 offers a natural-language task where headroom is uniformly high ($1-R^2_{\text{init}} = 0.69$--$0.90$ across embedding dimensions). If the headroom framework generalizes, we would predict a strong Frobenius--$\Delta$Error association here. This is confirmed: the config-mean Spearman correlation is $\rho=0.90$ ($p=0.037$, 5 dimensions $\times$ 5 seeds). The effective rank of the initial feature kernel is low ($p_{\text{eff}} \approx 3$), well below the test sample size ($n=2115$), which Proposition~\ref{prop:gap} predicts should yield stable diagnostic estimates. CKA is uniformly low ($0.01$--$0.03$)---attention-layer representations reorganize almost entirely during training regardless of scale, meaning the CKA blind spot (Regime~III) does not arise. A shuffled-label null baseline confirms that Frobenius change remains substantial ($145\%$ across configurations) even in the absence of signal, while $\Delta$Error becomes negative (performance degrades), consistent with the null analysis in Section~4.2. Together, the ViT and SST-2 results confirm that the dissociation is not an artifact of MLP parameterization or vision modalities, and that the headroom framework (Section~5.1.1) predicts diagnostic reliability across model families and data modalities.

To further test whether the headroom framework holds on architectures with batch normalization and skip connections, we trained a ResNet-18 on two tasks: binary MNIST (channel widths 16--48, 5 seeds each) and CIFAR-100 superclass binary (channel widths 32--96, 5 seeds each). The MNIST results are consistent with the headroom prediction: headroom strongly correlates with diagnostic reliability across widths ($\rho=1.0$, $p<0.001$, $n=4$ widths), though headroom is low ($0.03$--$0.08$) due to the task's ease. The CIFAR-100 results reveal a different scaling regime: unlike the synthetic experiments where wider networks change less and improve less, on this harder task wider networks actually learn more ($\Delta$Error increases from $0.09$ to $0.44$ with width) despite lower headroom ($0.84$--$0.70$). This inverts the headroom-$\Delta$Error correlation ($\rho=-1.0$), but the Frobenius-$\Delta$Error association remains strong ($\rho=0.80$). The effective rank is near-minimal ($p_{\text{eff}} \approx 1$) across all widths in both tasks, consistent with Proposition~\ref{prop:gap}'s prediction of stable diagnostic estimates. The CIFAR-100 null baseline confirms substantial Frobenius change under shuffled labels ($168\%$ on average). Together, these results confirm that the headroom framework generalizes across model families and task difficulties, though the specific relationship between width, headroom, and improvement is task-dependent and merits further study.

% ============================================================
\section{Towards a Reliability Map for Representation Diagnostics}
\label{sec:discussion}

We have applied RDEP to one widely used diagnostic (feature-kernel evolution magnitude) across representative regimes, producing three reliability regimes and a Diagnostic Validation Gap. The central lesson is that representation metrics themselves require predictive validation, just as models do.

\begin{figure}[t]
    \centering
    \includegraphics[width=0.8\textwidth]{fig5.png}
    \caption{\textbf{RDEP Reliability Map.} Observable kernel diagnostics enter three reliability regimes where common interpretations fail.  Regime~I: high CKA does not imply sufficient representation.  Regime~II: large Frobenius change does not imply useful task adaptation.  Regime~III: CKA $>0.99$ coexists with substantial spectral reorganization.  Below the regimes, a predictive framework (Section~5.1.1) shows that regime membership can be anticipated from headroom ($1-R^2_{\text{init}}$) and spectral alignment (in-sample $R^2=0.847$, LOOCV $R^2=0.710$), but instance-level corrections (Section~5.2) reveal a deeper validation gap ($\rho=-0.45$).}
    \label{fig:reliability_map}
\end{figure}

Figure~\ref{fig:reliability_map} summarizes the reliability regimes we have identified.  The general diagnostic bound (Section~2.5), which follows from existing theory~\citep{canatar2021spectral}, predicts that any label-invariant similarity diagnostic and risk change can in principle dissociate.  Our experiments measure how prevalent this dissociation is across representative regimes and identify specific conditions under which it occurs. These three regimes directly inform our validated diagnostic framework and validation principles (Sections~5.2--5.3).

\subsection{Implications for understanding feature learning}

Our results suggest that the relevant explanatory variable may not be the magnitude of kernel evolution, but how representations reorganize along task-relevant directions.  This perspective connects to the adaptive kernel framework~\citep{lauditi2025adaptive}, where feature-learning networks are described as kernel machines with data-dependent kernels: within that framework, kernel evolution is expected and necessary, but our results show it is not \emph{uniformly informative}.  The question shifts from ``how much did the kernel change?'' to ``what changed, and does it matter for the task?''  Identifying the appropriate measures of task-relevant representation change is an open question.  The loss landscape perspective~\citep{fort2019deep} offers a complementary lens: kernel stability at large widths may reflect the geometry of the loss basin rather than the absence of feature learning. Multiple explanatory frameworks converge on the same conclusion that kernel evolution magnitude alone is insufficient.

\subsubsection{Predicting Diagnostic Reliability from Pre-Training Quantities}

The stark task dependence we observe, $\rho=0.14$ (poly), $\rho=0.94$ (highfreq), $\rho=1.0$ (GMM), can be understood through two pre-trainable quantities: the target's alignment with the initial kernel's eigenspace, and the remaining room for improvement (headroom). To disentangle these, we constructed a family of target functions using Hermite polynomials of degrees 1 through 15, which smoothly vary in spectral complexity while remaining well-defined functions of the input. For each degree, we predicted the diagnostic reliability from (i) the concentration of the target projection $v_i$ across the top 3 eigendirections of the initial kernel (alignment) and (ii) $1-R^2_{\text{init}}$ from a linear probe on initial features (headroom). We then trained models across 6 widths and 5 seeds, measuring the realized Frobenius--$\Delta$Error Spearman $\rho$.

Across the 15 target functions, the dominant predictor of diagnostic reliability is headroom (Spearman $\rho=0.81$, $p=0.0002$): targets with large remaining room for improvement show stronger association between kernel change and generalization gain. Alignment alone shows no significant univariate correlation with reliability ($\rho=-0.19$, $p=0.49$), but interacts positively with headroom. A linear regression of $\rho$ on alignment, headroom, and their interaction achieves an in-sample $R^2=0.847$ (leave-one-out CV $R^2=0.710$) with a positive interaction coefficient, indicating that alignment amplifies the headroom effect when both are present.

This framework reinterprets the three original targets naturally. GMM has both high alignment (binary cluster structure projects onto dominant eigendirections) and high headroom (init R$^2\approx0$), yielding the strongest $\rho$. Poly has high alignment (init R$^2\approx0.68$) but low headroom. The initial representation already captures the target, leaving little room for kernel change to track improvement. Highfreq has low headroom at small widths (init R$^2\approx-0.13$ implies the initial representation is worse than the mean, so there is room to learn) but its low alignment means the observed kernel change occurs in task-orthogonal directions, explaining the width-confounded correlation structure. Unlike the earlier post-hoc narrative, the headroom-alignment hypothesis was specified before the Hermite experiment (the protocol and target functions were fixed prior to observing the outcome $\rho$ values). The same quantities (alignment and headroom) were computed from pre-training information alone and evaluated against the observed $\rho$ across all 15 targets. The regression and Spearman results confirm that diagnostic reliability is partially predictable from pre-training quantities, with headroom as the dominant factor.

It is important to distinguish two levels of prediction. Section~5.1.1 addresses a \emph{task-level} classification problem: given a new target, can we predict whether kernel-change magnitude will be a reliable diagnostic for that task? The answer is yes: a regression using headroom and alignment achieves in-sample $R^2=0.847$ (LOOCV $R^2=0.710$), with headroom as the dominant predictor (Spearman $\rho=0.81$, $p=0.0002$). Section~5.2 addresses a harder \emph{instance-level} regression problem: can we use the same spectral information to construct or correct an individual diagnostic so that it reliably tracks generalization for each specific model? This requires estimating $v_i$ and $\Delta\lambda_i$ from finite samples, which introduces a data-reuse problem that cross-validation exposes. The two results are not contradictory; they reveal a hierarchy of difficulty: regime classification is feasible, but instance-level correction faces a fundamental validation gap that remains open.

\subsection{Towards Task-Aware Kernel Diagnostics}

The dissociation we have documented suggests that the relevant quantity is not whether the kernel changes, but whether it changes in task-relevant directions.  A simple task-aligned projection shows this distinction. Projecting $\Delta K$ onto target-weighted eigendirections $\Delta_{\text{align}} = \sum_i v_i \cdot \Delta\lambda_i$ achieves a positive correlation with generalization improvement on the polynomial task when evaluated on the same data.  Cross-validating this measure (computing $v_i$ from 4/5 of the training data and evaluating improvement on the held-out 1/5, repeated over 5 folds and 5 seeds) yields a \emph{negative} correlation ($\rho=-0.45$, $p=0.02$, $n=25$).

This negative result reveals a fundamental \textbf{Diagnostic Validation Gap}. Estimating task alignment on the same data yields low bias but high optimism (overfitting). Cross-validation reduces optimism but introduces variance due to finite-sample eigendecomposition instability. This is directly analogous to the bias-variance trade-off in supervised learning. Our result ($\rho=-0.45$) is not a failure; it is the first empirical demonstration of this trade-off for representation diagnostics. Closing this gap is an open methodological question.

A practical alternative that avoids the validation gap is already available: the \textbf{Linear Probe Gain (LPG)}. It is defined as $\text{LPG} = R^2_{\text{post}} - R^2_{\text{init}}$, where $R^2_{\text{init}}$ and $R^2_{\text{post}}$ are linear probe coefficients of determination using initial and trained features (Supplementary Table~S1). Unlike $\Delta_{\text{align}}$, the linear probe uses a separate readout trained on held-in data and evaluated on held-out data. This avoids the data reuse problem entirely on logical grounds. We therefore recommend LPG as a \emph{principled} measure of task-relevant representation change that avoids data reuse.

We evaluate LPG under the same 5-fold cross-validation procedure used for $\Delta_{\text{align}}$ (polynomial task, width~64, 10 seeds). Its Spearman correlation with held-out gain is $\rho=-0.35$ ($p=0.01$, $n=50$), compared to $\rho=-0.42$ for $\Delta_{\text{align}}$. For comparison, the in-sample (non-cross-validated) LPG correlation with $\Delta$Error, computed across the full task--width configuration grid (3 datasets $\times$ 6 widths, configuration means), is $\rho=0.78$ ($p=0.0001$), with task-specific values of $\rho=0.77$ (poly) and $\rho=0.89$ (highfreq). LPG is less severely affected by the validation gap than spectral corrections, but it is not a panacea: both diagnostics suffer from the same fundamental limitation under finite samples and flat spectra. That LPG and $\Delta_{\text{align}}$ both yield negative cross-validated correlations is not a failure of either diagnostic in isolation, but evidence for an irreducible gap that Proposition~\ref{prop:gap} attributes to finite-sample estimation variance. The practical implication is that no existing diagnostic can be trusted without held-out validation---a meta-level conclusion that applies equally to our own proposed measure. Figure~\ref{fig:rdep} (Step~5) formalizes an adaptive selection rule that chooses between spectral and linear-probe diagnostics based on the effective rank of $K_0$, following the theoretical guidance of Proposition~\ref{prop:gap}.

\vspace{4pt}
\subsection{Principles for Diagnostic Evaluation}\label{sec:principles}
\vspace{2pt}

\subsubsection{RDEP in Practice: Preventing Diagnostic Misjudgments}

To illustrate how RDEP changes practical conclusions, Table~\ref{tab:rdeep_cases} walks through three representative configurations. It compares the interpretation a practitioner would draw without RDEP (relying on CKA or Frobenius alone) versus with RDEP (following the full workflow).

\begin{table}[t]
\centering\small
\caption{RDEP prevents three common diagnostic misjudgments.  For each configuration, the naive interpretation (based on CKA or $\Delta K$ alone) is contradicted by the RDEP workflow, which combines multiple metrics and width-controlled analysis.}
\label{tab:rdeep_cases}
\begin{tabular}{p{2cm}p{2.2cm}p{2.2cm}p{2.2cm}p{2.2cm}}
\toprule
Config. & Naive reading (CKA/$\Delta K$) & RDEP Step 1 (CKA) & RDEP Step 2 ($\Delta K$) & RDEP Step 3--4 (spectral + alignment) \\
\midrule
Poly w=32 & CKA=0.92, $\Delta K=25\%$ $\rightarrow$ moderate change, moderate learning? & CKA $\gg 0.99$? No. Cannot conclude ``no learning'' & $\Delta K >10\%$. Cannot conclude ``useful adaptation'' & $\rho=0.14$ (near-zero). Spectral alignment is weak. Correct conclusion: kernel change does not predict gain. \\
\midrule
Highfreq w=32 & $\Delta K=103\%$ $\rightarrow$ large change = useful learning? & CKA=0.85 $<0.99$. No false stability claim & $\Delta K >10\%$. Do not conclude adaptation & $\rho=0.94$ raw, but collapses to $\rho=0.10$ after width control. The apparent association is a width artifact. \\
\midrule
GMM w=128 & CKA=0.99, $\Delta K=53\%$ $\rightarrow$ contradiction? & CKA $\approx 0.99$. Check further & $\Delta K >10\%$. Do not conclude adaptation & $\rho=1.0$ even after width control ($\rho=0.91$). Alignment is genuine. Conclusion: for well-aligned tasks, $\Delta K$ \emph{can} be informative. \\
\bottomrule
\end{tabular}
\end{table}

Our findings translate into three principles for evaluating representation diagnostics, which parallel the reliability regimes identified in Section~4:

\begin{enumerate}
    \item \textbf{Global similarity $\neq$ representation sufficiency.}  High CKA between initial and trained kernels does not imply the absence of meaningful spectral reorganization.  CKA and spectral quantities carry distinct information, and both merit independent reporting.
    \item \textbf{Kernel magnitude $\neq$ task relevance.}  The magnitude of Frobenius change conflates task-aligned with task-orthogonal movement.  Evaluating whether representation changes are useful requires measuring task-relevant quantities such as target projections $v_i$ or linear probe performance.
    \item \textbf{Task alignment must be validated.}  Any diagnostic that uses the same data to estimate both target alignment and improvement risks overfitting (Section~5.2).  Predictive validation, through cross-validation or held-out tasks, is necessary before a diagnostic can be considered reliable.
\end{enumerate}

\begin{figure}[t]
\caption{RDEP Validation Protocol}
\label{fig:rdep}
\small
\textbf{Input:} Model family $\mathcal{F}$, task distribution $\mathcal{T}$, candidate diagnostic $\mathcal{D}$, significance level $\alpha$, number of widths $W$, number of seeds $S$ \\
\textbf{Output:} Reliability map $\mathcal{R}$ indicating regimes where $\mathcal{D}$ is informative about generalization

\begin{enumerate}
    \item \textbf{Calibration sweep.}
    For each task $f \in \mathcal{T}$, each width $w \in \{w_1, \dots, w_W\}$, and each seed $s \in \{1, \dots, S\}$:
    \begin{enumerate}
        \item Train model $M_{w,s} \in \mathcal{F}$ on $n$ samples from $f$.
        \item Compute initial and final kernels $K_0, K_T$, diagnostic $d_{w,s} = \mathcal{D}(K_0, K_T)$, and generalization improvement $\Delta e_{w,s}$.
    \end{enumerate}
    Construct calibration table $\mathcal{C} = \{\bar{d}_w, \overline{\Delta e}_w\}_{w=1}^W$ where $\bar{d}_w, \overline{\Delta e}_w$ are means over seeds.

    \item \textbf{Width-controlled analysis.}
    Compute partial Spearman correlation $\rho(\mathcal{D}, \Delta\mathcal{E} \mid \text{width})$ controlling for the width-driven trend.  If $|\rho_{\text{partial}}| < \alpha$, flag the association as a potential width artifact.

    \item \textbf{Spectral inspection.}
    Decompose $K_0 = \sum_i \lambda_i \phi_i \phi_i^\top$ and compute target projections $v_i = \langle \phi_i, y \rangle^2 / n$.  Report spectral concentration $c_k = \sum_{i=1}^k v_i / \sum_{i=1}^n v_i$ for $k \in \{3, 5, 10\}$.

    \item \textbf{Held-out validation.}
    If $\mathcal{D}$ incorporates estimated target projections $\hat{v}_i$, perform $K$-fold cross-validation:
    \begin{enumerate}
        \item For each fold, estimate $\hat{v}_i^{(\text{train})}$ on 4/5 of data, measure $\Delta\mathcal{E}$ on held-out 1/5.
        \item Compute $\rho_{\text{cv}}$ across all folds.  Report both in-sample and cross-validated correlations.
        \item If $\rho_{\text{cv}} \leq 0$ while $\rho_{\text{in}} > 0$, conclude that the diagnostic suffers from data-reuse overfitting (Diagnostic Validation Gap, Proposition~\ref{prop:gap}) and is not reliable without held-out validation.
    \end{enumerate}
    \item \textbf{Adaptive diagnostic selection.}
    Compute effective rank $p_{\text{eff}} = (\sum \lambda_i)^2 / \sum \lambda_i^2$ of $K_0$ and credibility score $\kappa = p_{\text{eff}} / n$:
    \begin{enumerate}
        \item If $\kappa < 0.3$ (concentrated spectrum): spectral diagnostics $\mathcal{D}_{\text{align}}$ are reliable; proceed with Step~4.
        \item If $\kappa > 0.7$ (flat spectrum): spectral diagnostics are dominated by estimation variance; fall back to LPG (Linear Probe Gain) as a more stable alternative.
        \item If $0.3 \leq \kappa \leq 0.7$: report both diagnostics and mark the regime as ambiguous.
    \end{enumerate}
\end{enumerate}
\end{figure}

These principles suggest a practical workflow:
\begin{center}
\small
\begin{tabular}{p{0.92\textwidth}}
\emph{Step 1:} Compute CKA$(K_0, K_T)$. If CKA $> 0.99$, do \emph{not} conclude ``no feature learning.'' \\
\emph{Step 2:} Check Frobenius change ($\Delta K$). If $\Delta K > 10\%$, do \emph{not} conclude ``useful adaptation'' unless the task-kernel alignment is independently confirmed. \\
\emph{Step 3:} Examine spectral quantities (eigenvalue ratios, target projections $v_i$). \\
\emph{Step 4:} Only after Step 3, assess whether representation changes are task-aligned.
\end{tabular}
\end{center}
This workflow implements the three principles: it first prunes the common misinterpretations (Principles~1 and~2), then requires spectral inspection and task-aligned validation (Principle~3) before drawing conclusions about feature learning.

\subsection{Limitations}
This work studies \emph{diagnostic validity}. That is, we ask whether observable kernel evolution provides a reliable signal about generalization, rather than studying causal mechanisms.  We do not claim that kernels are irrelevant to generalization; NTK theory remains valuable in its intended regime, and similarity metrics capture aspects of representation geometry.  Our contribution is bounding the reliability of these metrics as diagnostics.

Specific limitations:
\begin{itemize}
    \item \textbf{Architecture scope.}  Our main experiments use two-layer MLPs for tractability and direct connection to NTK theory.  Section~4.4 extends the analysis to CNNs, Vision Transformers, NLP Transformers, and ResNets, confirming that the dissociation is not an artifact of two-layer architectures.  However, the CIFAR-100 results reveal a different scaling regime (wider networks learn more despite lower headroom), suggesting that the headroom framework's predictions are task-dependent in ways not yet fully characterized.  Larger-scale architectures (e.g., ResNet-50 on ImageNet, LLMs) remain open challenges.  While the headroom framework's consistency across five distinct architectures suggests general validity, this is not yet proven at scale.
    \item \textbf{Feature kernel, not NTK.}  Whether analogous results hold for Jacobian-based kernels is open.
    \item \textbf{Standard parameterization confirmed in NTK param.}  We verified the dissociation on the polynomial task under NTK parameterization (init variance $1/d$). The pattern is qualitatively identical: at width~32, Frobenius change is $38\%$ with gain $-0.036\pm0.038$, and at width~128, Frobenius drops to $8\%$ with gain $+0.022$. The dissociation does not arise from a specific parameterization choice.
    \item \textbf{Highfreq results are seed-dependent, not a capacity artifact.}  With 10 seeds at width~32, the highfreq target shows mean gain $0.190\pm0.130$. Kernel evolution at small widths does capture \emph{some} task-relevant information on average, but the standard deviation means individual seeds range from near-zero to substantial improvement.  The original single-seed result (gain$\approx0$) was not representative.  A controlled comparison with the low-frequency target $\sin(x_1)+\cos(x_2)$ at the same width shows \emph{negative} gain ($-0.037\pm0.042$). This confirms that target structure, not network capacity, determines whether kernel evolution helps or hurts.
    \item \textbf{Synthetic data.}  Real-world tasks may present additional structure.
    \item \textbf{Hermite regression sample size.}  The headroom-alignment regression in Section~5.1.1 is based on 15 Hermite target functions, which still limits the precision of the leave-one-out CV estimate ($R^2=0.710$).  A larger target family would provide a tighter bound on held-out predictive accuracy.
    \item \textbf{RDEP scope.}  The RDEP framework requires spectral decompositions, which become computationally expensive for very large models (e.g., Transformers with hundreds of millions of parameters).  Our analysis focuses on feature kernels from two-layer MLPs and small CNNs, with extensions to small Vision Transformers and NLP Transformers (Section~4.4).  The positive SST-2 result ($\rho=0.90$) provides cross-modal evidence that the headroom framework generalizes beyond vision, but the calibration patterns for large-scale attention architectures or Jacobian-based (NTK) kernels remain open.  Extending to large-scale Transformers and LLMs is computationally prohibitive but conceptually straightforward; we leave this to future work.
\end{itemize}

\subsection{Looking Forward}\label{sec:future}

The framework developed here, calibrating diagnostic reliability through controlled perturbation, extends beyond feature-kernel evolution.  Any representation similarity metric (CKA, SVCCA, PWCCA, RSA, CCA) is label-invariant and therefore subject to the same bound (Section~2.5).  Our empirical methodology, systematically varying task structure, width, and seed while measuring the metric-improvement correlation, provides a template for evaluating the reliability of these metrics in turn.  When spectral decompositions are computationally expensive, we recommend Linear Probe Gain (LPG) as a practical RDEP-compatible alternative that avoids the data-reuse pitfall of spectral task-aligned measures.

Beyond feature-kernel evolution, this work establishes a broader principle: representation diagnostics should be held to the same evidential standard as the models they evaluate. Just as we demand cross-validation for predictors, we must demand predictive validation for the tools we use to interpret them. RDEP is a step toward that standard, but the open challenge is a community-wide shift in how we evaluate diagnostic tools, from CKA to saliency maps to feature attribution.

To facilitate standardized evaluation, we define \textbf{DiaBench}, a diagnostic benchmark built from the Hermite target family (Section~5.1.1) and the three calibration regimes (Section~4.1--4.3). DiaBench provides a controlled testbed where the ground-truth reliability of any candidate diagnostic can be measured against held-out data across a known spectrum of task-model alignment and headroom. The benchmark spans three difficulty levels: Regime~I (stability without learning), Regime~II (change without predictable learning), and Regime~III (CKA blind spot). Accompanying code is planned for release as an open-source library (\texttt{rdep-tools}) implementing the RDEP protocol, DiaBench data generators, and cross-validation modules for any similarity metric. We hope this infrastructure enables the community to empirically validate---rather than assume---the reliability of representation diagnostics. Future work should treat representation diagnostics as predictive tools requiring validation through held-out tasks, cross-validated projections, and explicit characterization of failure regimes.

% ============================================================
\bibliographystyle{tmlr}
\bibliography{refs}


% ============================================================
\appendix
\setcounter{figure}{0}
\setcounter{table}{0}
\renewcommand{\thefigure}{S\arabic{figure}}
\renewcommand{\thetable}{S\arabic{table}}
\section{Supplementary Analysis and Figures}

\subsection*{Task-Aligned Kernel Change}
\label{sec:task-aligned}
The task-aligned measure $\Delta_{\text{align}} = \sum_i v_i \cdot \Delta\lambda_i$ projects $\Delta K$ onto target-weighted eigendirections. On the polynomial task, the task-aligned and raw measures give qualitatively different correlations, though both are noisy at this sample size. On highfreq, the correlation remains near zero, consistent with the finding that substantial kernel evolution can occur in task-orthogonal directions without systematic generalization improvement (Section~5.1). We caution that $v_i$ and the improvement are computed on the same data; Section~5.2 shows that cross-validation flips the correlation to negative ($\rho=-0.45$), reinforcing that task-aligned measures must be validated on held-out data.

\section{Supplementary Figures}

\begin{figure}[ht]
    \centering
    \includegraphics[width=0.8\textwidth]{figS1.png}
    \caption{Spectral redistribution. For each dataset at width 128, top panels show eigenvalue ratio $\lambda_{\text{post}}/\lambda_{\text{init}}$ per eigen-direction; bottom panels show target projections $v_{i}$ before and after training. GMM exhibits substantial reorganization in both quantities despite high CKA.}
    \label{fig:spec_redist}
\end{figure}

\begin{figure}[ht]
    \centering
    \includegraphics[width=0.6\textwidth]{mnist_scatter.png}
    \caption{MNIST dissociation. Kernel change versus generalization improvement for MNIST binary classification (0 vs.\ 1).  The association is moderate (Spearman $\rho=0.66$, $p<0.001$), indicating that the dissociation is less extreme on real data: kernel change carries \emph{some} predictive signal but remains incomplete as a sufficient diagnostic.}
    \label{fig:mnist_scatter}
\end{figure}

\begin{figure}[ht]
    \centering
    \includegraphics[width=0.65\textwidth]{figS2.png}
    \caption{CKA versus SVCCA across all configurations (10 seeds per condition).  Points are colored by dataset.  The two metrics show strong agreement (pooled Spearman $\rho=0.89$, $p<0.001$), confirming that the dissociation documented with CKA is not specific to one metric choice.}
    \label{fig:cka_vs_svcca}
\end{figure}

\begin{figure}[ht]
    \centering
    \includegraphics[width=0.6\textwidth]{fig_epoch_dynamics.png}
    \caption{Training dynamics for poly, width 32 (illustrative single-seed trajectory). Test error drops while CKA (right axis) stabilizes within 50 epochs. Across 10 seeds, the mean improvement is $0.003\pm0.049$ (Table~\ref{tab:results}), so this trajectory shows the temporal decoupling mechanism rather than the typical effect size.}
    \label{fig:epoch}
\end{figure}

\begin{figure}[ht]
    \centering
    \includegraphics[width=0.85\textwidth]{figS5.png}
    \caption{CNN eigenvalue spectrum evolution on CIFAR-10 before (blue) and after (orange) training, for channel widths $W=16$ (left) and $W=128$ (right).  At $W=128$, the top eigenvalues increase substantially after training while the rank order is preserved, confirming that spectral reorganization persists even when CKA is high ($0.991$).}
    \label{fig:cnn_spectrum}
\end{figure}

\begin{table}[ht]
    \centering\footnotesize
    \setlength{\tabcolsep}{4pt}
    \caption{CIFAR-10 CNN results (5 seeds per width).  Values are means $\pm$ s.d.  The linear probe $\Delta R^2$ peaks at $W=32$ ($+0.519\pm0.306$) while Frobenius change shows a non-monotonic pattern.}
    \label{tab:cnn_cifar}
    \begin{tabular}{lccc}
    \toprule
    Width & CKA & $\Delta K$ (\%) & $\Delta R^2$ \\
    \midrule
    16  & $0.919\pm0.039$ & $55.1\pm9.8$  & $+0.400\pm0.246$ \\
    32  & $0.947\pm0.024$ & $62.5\pm19.4$ & $+0.519\pm0.306$ \\
    64  & $0.962\pm0.027$ & $65.9\pm20.5$ & $+0.509\pm0.321$ \\
    128 & $0.991\pm0.007$ & $43.4\pm11.9$ & $+0.149\pm0.125$ \\
    256 & $0.989\pm0.009$ & $51.3\pm16.9$ & $+0.266\pm0.217$ \\
    \bottomrule
    \end{tabular}
\end{table}

\begin{table}[ht]
    \centering
    \caption{MNIST binary classification (0 vs.\ 1).  Width-dependent kernel stability shows a non-monotonic pattern: CKA rises from 0.768 to 0.963 as width increases from 32 to 256, then \emph{declines} at width 1024 ($0.867\pm0.024$, consistent across all 5 seeds).  The Frobenius change similarly increases at width 1024 ($81.5\pm9.3\%$), suggesting that on high-dimensional inputs (784D) very wide networks enter a distinct optimization regime where the kernel continues to reorganize substantially.  The cause of this reversal is not yet established but appears robust across seeds.  Despite these variations, the dissociation between $\Delta K$ and generalization improvement persists (Spearman $\rho=0.66$, $p<0.001$; see Supplementary Figure~S2).}
    \label{tab:mnist}
    \begin{tabular}{lccc}
    \toprule
    Width & CKA & $\Delta K$ (\%) \\
    \midrule
    32   & $0.768\pm0.090$ & $237.4\pm129.8$ \\
    64   & $0.877\pm0.030$ & $71.7\pm33.5$   \\
    128  & $0.917\pm0.026$ & $39.1\pm7.0$    \\
    256  & $0.963\pm0.011$ & $50.3\pm10.7$   \\
    512  & $0.963\pm0.028$ & $51.8\pm16.5$   \\
    1024 & $0.867\pm0.024$ & $81.5\pm9.3$    \\
    \bottomrule
    \end{tabular}
\end{table}

\begin{table}[htbp]
    \centering\footnotesize
    \setlength{\tabcolsep}{4pt}
    \caption{Linear probe R$^2$ on frozen features (mean over 5 seeds).  Poly features improve most when kernel change is large.  Highfreq features degrade after training despite $>90\%$ kernel change.}
    \label{tab:linear_probe}
    \begin{tabular}{lcccc}
    \toprule
    Dataset & Width & Init R$^2$ & Post R$^2$ & $\Delta$R$^2$ \\
    \midrule
    poly & 32  & $0.682$ & $0.961$ & $+0.279$ \\
    poly & 64  & $0.870$ & $0.963$ & $+0.093$ \\
    poly & 128 & $0.953$ & $0.970$ & $+0.017$ \\
    poly & 256 & $0.969$ & $0.975$ & $+0.006$ \\
    poly & 512 & $0.984$ & $0.985$ & $+0.001$ \\
    poly & 1024 & $0.987$ & $0.988$ & $+0.000$ \\
    \midrule
    highfreq & 32  & $-0.126$ & $-0.867$ & $-0.741$ \\
    highfreq & 64  & $-0.056$ & $-0.611$ & $-0.555$ \\
    highfreq & 128 & $-0.284$ & $-0.632$ & $-0.348$ \\
    highfreq & 256 & $-0.321$ & $-0.518$ & $-0.197$ \\
    highfreq & 512 & $-0.547$ & $-0.536$ & $+0.011$ \\
    highfreq & 1024 & $-0.515$ & $-0.468$ & $+0.047$ \\
    \bottomrule
    \end{tabular}
\end{table}

\begin{table}[h!]  % force immediate placement
    \centering\footnotesize
    \setlength{\tabcolsep}{4pt}
    \caption{FashionMNIST binary classification.  The Spearman correlation between kernel change and generalization improvement is \emph{negative} (Spearman $\rho=-0.43$, $p=0.017$).  See Section~4.4 for discussion.}
    \label{fig:fashion}
    \begin{tabular}{lccc}
    \toprule
    Width & CKA & $\Delta K$ (\%) & $\Delta$Error \\
    \midrule
    32   & $0.510\pm0.042$ & $67.0\pm5.2$  & $+8.670$ \\
    64   & $0.587\pm0.038$ & $77.2\pm4.1$  & $+1.355$ \\
    128  & $0.667\pm0.031$ & $84.7\pm3.8$  & $-0.032$ \\
    256  & $0.725\pm0.028$ & $91.6\pm2.9$  & $+0.010$ \\
    512  & $0.763\pm0.025$ & $94.2\pm3.1$  & $+0.108$ \\
    1024 & $0.747\pm0.030$ & $97.3\pm2.6$  & $+0.115$ \\
    \bottomrule
    \end{tabular}
\end{table}

\begin{table}[h!]  % force immediate placement
    \centering\footnotesize
    \setlength{\tabcolsep}{4pt}
    \caption{Random-label baseline (poly target, shuffled labels).  Kernel change remains substantial at small widths despite the absence of any signal, demonstrating that large Frobenius change alone cannot distinguish signal-driven adaptation from noise.  Results are means over 5 seeds.}
    \label{tab:random_label}
    \begin{tabular}{lccc}
    \toprule
    Width & CKA & $\Delta K$ (\%) & $\Delta$Error \\
    \midrule
    32   & $0.851\pm0.016$ & $103.0\pm7.3$   & $+0.111$ \\
    64   & $0.944\pm0.014$ & $49.2\pm3.1$    & $+0.028$ \\
    128  & $0.984\pm0.002$ & $21.2\pm1.6$    & $-0.003$ \\
    256  & $0.997\pm0.000$ & $6.8\pm1.9$     & $+0.004$ \\
    512  & $1.000\pm0.000$ & $2.1\pm0.3$     & $+0.001$ \\
    1024 & $1.000\pm0.000$ & $1.0\pm0.2$     & $+0.000$ \\
    \bottomrule
    \end{tabular}
\end{table}

\end{document}
